\chapter{Project Commentary}
When starting my EPQ I set out with the question: How can we improve our ability to study information? This was a question I had long been interested in as I progressed through the education system. To help me refine my project idea I looked at the variety of steps I found myself going through during the process of learning and zoomed in on the area of note-taking. As the first step any student will take in recording information during classes, I had touched on various note-taking methods in the past such as Cornell notes, yet I still found myself assessing many half completed, roughly written pieces of paper and digital folders full of documents. I was certain the information I needed was there somewhere, but precisely where, I had no idea.

This frustration with past disorganisation connected with an area of interest in which I had been building my skills: computer science. In the past I have built a variety of small tools such as a music playing website and advertising pages for local companies. I found these processes enjoyable as I shaped my vision into a tangible product, so much so I now envisioned myself working within this industry as a career. 

Bridging these two goals together then became the core focus of my project: to build my own platform that I could use to improve my learning, designed with psychological research at the very centre. I began my project by creating a rough Eisenhower matrix containing the various features I would want my platform to have. This was crucial in matching the scope of my software development with the skills required and resources I had available to maximise the impact of my end artefact.

When beginning to research the problem of note-taking, my first step was to survey a wide range of participants from students, teaching faculty and those in corporate careers to name but a few. My project's primary objective was to address the needs that people have in a wide range of situations, and so I wanted to keep their wants and needs at the centre of my research. From the survey responses received and verbal feedback three themes emerged:

\begin{enumerate}
    \item A desire to have easy access to all their notes anywhere
    \item A need for much better systems to store ideas and crucially find them later
    \item The ability to seamlessly combine typed notes with quick scribbles and diagrams  on a touch screen device. 
\end{enumerate}

With these ideas in mind I then set out to analyse how competitors in the market currently tackle these challenges. Through my primary research I identified numerous common options such as Notion, Obsidian and Evernote; more niche options such as Capacities and crucially, default options present on many users mobile devices such as Apple notes, being the most commonly used platforms thanks to their ease of use and low bar to access. This research of existing platforms served highly useful as it jump-started my understanding of the market allowing me to benefit from the lessons all of these platforms had faced before me.

Whilst I was conducting this primary research I was simultaneously learning from a vast array of sources about methods we can use to improve our learning techniques. Whilst users already knew what didn't work for them (having never been exposed to many of these learning techniques) I knew it was crucial to introduce methods which could work, but that they had never tried.

One technique that immediately stood out to me was Mind Maps. Traditionally a very free flowing and flexible tool used for drafting ideas together I found through my research that they were hugely beneficial in making neurological links between ideas that allows our minds to associate and therefore more easily jump to the required information as it follows a thought strand. This research into Mind Maps then led me into exploring another source about a developer's attempt at making a file management system called TagStudio (I'm a Mess, so I'm Making My Own File Organizer [TagStudio], 2024). Whilst I approached this source hoping to learn more about the technical details of his implementation, I soon found myself fascinated by the theory of using tags to link files rather than nested folders.

Unlike a traditional folder-based system (as almost every platform provides) a true tag based system in my context allows notes to be associated with any given tag, and crucially tags to also have other tags applied to them. This enables many powerful features such as filtering notes based on tags that they may (or may not) contain to find exactly the content you want. To illustrate this point imagine a document called “Maths homework 2025”, this document could have the tags: 2025, Maths, and homework all applied. The Maths tag could then have a “schoolwork” tag applied and suddenly by searching for any of these 4 different ideas the target document can be found, much more flexible than a traditional folder-based workflow

These two principles: Mind Maps and note tagging then formed the core philosophy of NeuroNote. By rethinking the way users fundamentally structured their notes I believed that I could dramatically improve recollection usability. Whilst the philosophy of idea tagging was strong its common flaw was a difficulty in understanding the process and how ideas are linked, My idea was for users to organically interact with a mind map, adding tags between notes to form visual links between them. These tags themselves would appear as 'nodes' on the mind map thus making the entire system very visual and searchable. Users wouldn't have to consider the complexity of naming and assigning tags, just dragging links between ideas and the system would automatically handle the rest.

Whilst ideologically sound, I didn't just want to rest on my academic research with these theories however so I soon set to work building a prototype system that I could present to some sample users. 

Throughout my development journey, prototypes were critical to help me quickly relay my findings to sample users, before investing a lot of time on an idea that was strong on paper but would face major issues when applied in practice. 

For this prototype I researched a technical 'library' named Cytoscope Js. This tool was initially designed for biologists to link together subcellular structures however I saw potential in the methods it used to relay information. With Cytoscape and a lot of custom code supporting it, I soon had a prototype to show users and they loved it. Whilst still very rough around the edges it proved the core idea feasible and incredibly beneficial to my participants note-taking journey, and the feedback was very positive. From this prototype I also gained some very useful advice. One such example being a participant noting it could be challenging see which notes were related to any other given note on a very dense graph. To solve this, I implemented a selection system in which clicking on a note or tag within the graph would make the selected object have an orange hue, whilst 'children' (notes which had the selected node as a parent either directly or through intermediates) gain progressively lighter highlights as they became more distant relations.

This process of Research, prototype, evaluate continued throughout my project as every core feature utilised a combination of Academic research, end user consultation, prototype development and evaluation to continually refine my artefact. 

An additional example of my process can be seen in the development of my whiteboard editor. I drew on Mueller and Oppenheimer's (2014) study “The Pen is Mightier than the Keyboard”, which argues that handwriting information improves cognitive performance as the physical act of notetaking encourages deeper processing of material. This finding, supported by further studies I reviewed aligned with feedback from my users who valued the flexibility of handwriting alongside typed notes.
To address this, I designed a specification using an Eisenhower matrix and began building an improved 'whiteboard' space for handwritten note-taking. This proved to be one of the most technically challenging aspects of the project, as I needed to support input across devices with very different mechanisms (mouse, trackpad, touchscreen). Overcoming these challenges helped me develop skills in troubleshooting and user interface design, particularly in communicating features in a simple and intuitive way.
After seeking feedback from users, they loved the feature-sets of my whiteboard and small 'quality of life' add-ons such as being able to insert pan-able whiteboard preview boxes within text-based documents (which additionally create links between the document and the whiteboard in the graph view thanks to my tagging system).

My final Artefact software was drawn together from all these prototypes in combination with a vast array of other skills such as GIT version control, web authentication \& security research and Server hosting with database management to name but a few. 

When I began this project, I set out with the broad question: How can we improve our ability to study information? Through my research, development, and user testing, I believe NeuroNote provides a meaningful step towards that goal. By combining mindmap-style visual links with a flexible tag-based system, my platform tackles the problems of disorganisation, accessibility, and knowledge recall that both my survey participants and I experienced. While my artefact is not without limitations, such as the complexity of very large graphs and the need for further optimisation across devices, the project has shown that integrating psychological research with user-centred software design can make studying more efficient and intuitive. Most importantly, this project has answered my question in a practical way: we can improve our ability to study information by rethinking how notes are captured, stored, and connected, making the learning process both more accessible and more engaging. Alongside these academic outcomes, the project has also developed my technical and transferable skills and confirmed my ambition to pursue computer science as a career.

For further evaluation of my complete project (including commentary on my user feedback) please visit my presentation evaluation above in addition to my Evaluation chapter. Additionally, for a much more in depth look at the process of researching and creating my artefact please visit my Research and Diary chapters below.
